Este trabalho apresentou uma extensão do modelo de tempo de promoção para a modelagem da componente de cura
por meio da incorporação de funções kernel, resultando no modelo GG-KM. Essa abordagem permite maior flexibilidade
na estimação, tornando possível capturar relações complexas e não lineares entre as covariáveis e a fração
de cura.

Nos experimentos com dados simulados, os resultados evidenciaram que kernels não lineares apresentaram
desempenho superior em cenários com maior complexidade, particularmente nos métodos 2 e 3, nos
quais alcançaram os menores valores de IBS. Em contrapartida, o kernel linear apresentou desempenho inferior
nesses cenários, reforçando a importância da flexibilidade introduzida pelos kernels.

Além disso, tanto nos dados simulados quanto na aplicação ao conjunto de dados reais, observou-se bom ajuste
entre a curva de sobrevivência estimada pelo modelo GG-KM e o estimador de Kaplan-Meier. Esse comportamento
fornece evidências de que o modelo proposto é capaz de representar adequadamente a fração de cura presente
nos dados.

De forma geral, os resultados obtidos indicam que o GG-KM constitui uma alternativa promissora para modelagem
de sobrevivência com fração de cura, especialmente em contextos onde a estrutura da relação entre covariáveis
e sobrevivência apresenta maior complexidade.
